\section{The failure that doesn't announce itself}
\label{sec:8.1}
A refusal announces itself: someone asked for an act and was turned down, and the person who asked can complain. A wrongful compliance looks like any other completed task. The person harmed usually isn't the requester and may never learn that a system named, scored, located or classified them.

An appeal route can't find that case, since it starts from a refusal, and reviewing every compliance would mean re-adjudicating everything. Three remedies follow: notice to the people affected, sampled re-adjudication by reviewers who didn't produce the record, and split custody of the logs.

The notice requirement: where a compliance names, scores, locates or classifies an identifiable person who didn't make the request, that person is notified, on the model of consumer credit's adverse-action notice, and has a right to respond and a right to contest. All three are unconditional, and none is exchanged for testimony: nobody has to give an account of their loss to be told what was done or to contest it. \emph{TransUnion v. Ramirez}, a Fair Credit Reporting Act case, held that a statutory violation without concrete harm gives no standing to sue in federal court \autocite{transunion2021}. That limits one thing, an individual suit for damages over a missed notice, which attaches where the compliance was acted on against the person: a denial, a detention, a place on a kill list. It doesn't limit notice, response or contest. Nor does it limit the standing to sue of organizations that represent a class of people affected. Where the person affected can't be identified (a target in a classified pipeline), individual notice doesn't bind, but notice to the class does, and that fact is published beside the list of positions where a refusal can be defeated.

Someone the floor protects, who thinks it's set wrongly or is protecting them in a way they'd decline, uses the same route. Notice tells them the floor acted on them, and contest reaches the floor court.
